\section{Introduction}\label{sec:intro}

Core decomposition ranks vertices by the amount of cohesive structure that still surrounds them after weak vertices are removed~\cite{kcoreSeidman,coreAlg,communityFinder}.
%
The classical \(k\)-core uses edges as the witness of cohesion.
%
Many graph mining tasks need a stronger witness, because a group may pass an edge-degree threshold while containing few triangles, \(4\)-cliques, or larger coordinated patterns.
%
The \emph{\((1,s)\)-nucleus decomposition}~\cite{nucleusSariyuce14} gives this vertex-centric higher-order ranking.
%
For a fixed \(s\), it assigns each vertex the largest value \(k\) for which the vertex belongs to a maximal region where every vertex participates in at least \(k\) internal \(s\)-cliques.
%
At \(s=2\), this is the \(k\)-core.
%
At \(s=3\), it ranks vertices by triangle-supported cohesion.
%
At larger \(s\), it exposes progressively tighter clique-based kernels.

Exact high-\(s\) decomposition is hard because the number of \(s\)-cliques can be enormous.
%
The current state of the art avoids explicit enumeration with the Clique Path Index, or \cpi{}~\cite{NuclearCD}.
%
A \cpi{} path stores a hold set and a pivot set.
%
It compactly represents all \(s\)-cliques formed by taking every hold vertex and choosing the remaining vertices from the pivot set.
%
This representation made general \((r,s)\)-nucleus decomposition practical, but the general algorithm treats the index as mutable residual state during peeling.
%
It stores a mutable recursion tree, per-path vertex sets, reverse maps, and a priority queue.
%
This generality is costly.
%
It is also necessary for the general case, because when \(r\ge2\), peeling can delete edges or higher-order objects and can break the clique structure that a path stores.

The vertex-centric case has a stronger invariant.
%
When \(r=1\), peeling removes vertices but does not delete edges.
%
Therefore, a set of vertices that formed a clique when the \cpi{} was built remains a clique throughout the peel.
%
The path may stop contributing because a required hold vertex has been removed.
%
Or its contribution may shrink because one pivot vertex has been removed.
%
But the stored hold and pivot sets do not need to be rewritten.
%
The dynamic state of a path is only a liveness bit, its fixed target, and an effective pivot counter.

This observation is the thesis of the paper:
for \((1,s)\)-nucleus decomposition, vertex peeling preserves the \cpi{} structure, so exact decomposition can be performed over a static clique-path index using only counter updates.
%
Theorem~\ref{thm:vertex-removal} proves the static-\cpi{} invariant.
%
Lemma~\ref{lem:support-delta} turns each peel event into a closed-form binomial support delta.
%
Together they replace mutable residual-index maintenance with arithmetic over a static incidence layout.

Once the index is static, the core computation can be viewed through the usual \(h\)-index characterization of core values~\cite{kcoreLocal,lu2016hcore,hindex}.
%
\localH{} is the direct synchronous formulation.
%
It repeatedly refreshes every vertex value from the \cpi{} witness counts until the \(h\)-values stop changing.
%
\peelH{} is the event-driven realization of the same computation.
%
It finalizes vertices in nondecreasing value order and refreshes only vertices whose path counters changed.
%
The two algorithms compute the same fixed point, but \peelH{} avoids the global redundancy of synchronous rounds.

Making peeling cheap changes the bottleneck.
%
After support maintenance no longer mutates the residual \cpi{}, construction becomes a dominant part of end-to-end time on many benchmark instances.
%
This motivates \paraBuild{} as part of the same chain rather than as a separate optimization.
%
The degeneracy-ordered seed loop already gives disjoint ownership of \(s\)-cliques.
%
Because the downstream consumer needs only a static path list and incidence arrays, each thread can stream emitted paths into a local arena and merge the arenas deterministically.

The same deletion order also gives a core-value hierarchy.
%
\peelH{} emits vertices in nondecreasing \(\kappa_s\), because its bucket cursor never retreats.
%
Let \(V_k=\{v:\kappa_s(v)\ge k\}\).
%
The edge-connected components of \(G[V_k]\) can therefore be recovered by scanning this deletion log in reverse.
%
\buildHier{} uses only the input graph and union-find to extract this elder-rule join tree.

The contributions are a dependency chain.
\begin{itemize}[leftmargin=1.3em,itemsep=1pt,topsep=2pt]
\item \textbf{Static-\cpi{} state for \((1,s)\).} We prove that \cpi{} paths support exact vertex peeling without modifying their stored hold and pivot sets; every event becomes a path-death or pivot-counter update with a closed-form support delta (Theorem~\ref{thm:vertex-removal}, Lemma~\ref{lem:support-delta}).
\item \textbf{Two fixed-point realizations over the static index.} We formulate \localH{} as the synchronous \(h\)-index iteration on \cpi{} witness counts, and \peelH{} as the event-driven peel that computes the same fixed point without repeated global refreshes (Algorithms~\ref{alg:localH} and~\ref{alg:peelH}, Lemma~\ref{lem:equivalence}).
\item \textbf{Bottleneck-shifted static construction.} Once peeling becomes counter-only, \cpi{} construction becomes the dominant remaining cost on many instances; \paraBuild{} exploits the static output contract by streaming paths into thread-local arenas and merging them deterministically without hot-path synchronization (Algorithm~\ref{alg:parabuild}).
\item \textbf{Core-value hierarchy from the deletion log.} We show that the sorted log produced by \peelH{} is enough to recover the edge-connected super-level join tree in \(O((n+m)\alpha(n))\) time using only the input graph (Algorithm~\ref{alg:buildhier}).
\item \textbf{Experimental and application evidence.} We compare against the mutable-\cpi{} baseline, test output agreement, report runtime and memory, measure phase behavior and parallel construction, and use case studies to evaluate whether the \(r=1\) specialization is practically useful.
% TODO: verify the headline ten-graph, 762-configuration speedup and memory claims against the local evidence files in figures/ and the external data referenced by the plotting scripts.
\end{itemize}

The rest of the paper follows this chain.
%
Section~\ref{sec:prelim} defines the problem, the \cpi{}, and the mutable baseline.
%
Section~\ref{sec:static} proves the static-\cpi{} invariant.
%
Section~\ref{sec:hindex} derives \localH{} and \peelH{} from the same \(h\)-index view.
%
Section~\ref{sec:parabuild} parallelizes construction after peeling becomes cheap.
%
Section~\ref{sec:buildhier} extracts the hierarchy from the deletion log.
%
Sections~\ref{sec:exp} and~\ref{sec:casestudy} report performance and application evidence.
%
Sections~\ref{sec:related} and~\ref{sec:conclusion} position and conclude the paper.
